\chapter{Derivada Direccional y Valores Extremos de una Función de Varias Variables}

El cálculo diferencial de una variable responde a una pregunta geométrica
sencilla: ¿cuál es la pendiente de la recta tangente a una curva en un
punto? Cuando la curva está en el plano, hay una sola dirección posible
para la tangente. Pero en el espacio tridimensional, la situación cambia
radicalmente: desde un punto sobre una superficie parten infinitas
direcciones, y la pendiente de la recta tangente a la superficie depende
de cuál de esas direcciones se elija. Las derivadas parciales
$f_x$ y $f_y$ son casos particulares: miden la pendiente de la
superficie cuando nos movemos paralelamente a los ejes $x$ e $y$,
respectivamente. Pero para comprender plenamente el comportamiento de una
función de varias variables, se necesita una herramienta que mida su tasa
de cambio en \emph{cualquier} dirección del espacio. Esa herramienta es
la \textbf{derivada direccional}, y el vector que organiza toda esta
información es el \textbf{gradiente}.

La derivada direccional responde a la pregunta: si desde un punto me
desplazo una distancia infinitesimal en una dirección dada, ¿cuánto
cambia el valor de la función por unidad de longitud recorrida? El
gradiente, por su parte, no solo permite calcular la derivada direccional
en cualquier dirección mediante un producto escalar, sino que además
señala la dirección de máximo crecimiento de la función.

Esta lección está organizada en dos grandes bloques. En el primero se
introduce el gradiente y la derivada direccional, se establece la fórmula
que los relaciona y se estudia cómo maximizar la derivada direccional.
Además, se extiende el concepto de plano tangente a superficies dadas en
forma implícita mediante el gradiente. En el segundo bloque se aborda la
búsqueda de valores extremos (máximos y mínimos) de funciones de varias
variables, incluyendo el criterio de la segunda derivada para
clasificar puntos críticos y el Teorema del Valor Extremo para funciones
continuas sobre conjuntos cerrados y acotados.

%%==========================================================
\section{Gradiente y derivada direccional de una función multivariable}
%%==========================================================

Recordemos que las derivadas parciales de una función $f(x,y)$,
denotadas $f_x(x,y)$ y $f_y(x,y)$, indican cómo cambia la función
cuando nos movemos a lo largo de los ejes $x$ e $y$, respectivamente.
Es decir, muestran la pendiente de la recta tangente a la superficie
$z=f(x,y)$ en esas direcciones específicas. El objetivo de esta sección
es extender esta idea y aprender a medir cómo cambia la función en
\emph{cualquier} dirección, no solo a lo largo de los ejes. Esta medida
se llama \textbf{derivada direccional}. Antes de definirla,
necesitamos conocer un concepto clave: el \textbf{vector gradiente}.

%%----------------------------------------------------------
\subsection{Gradiente}
%%----------------------------------------------------------

\begin{definition}[Gradiente de una función escalar]
Sea $f$ una función escalar de dos o tres variables.
\begin{itemize}
    \item Si $f(x,y)$ es una función de dos variables, el \textbf{gradiente}
    de $f$ es el campo vectorial $\nabla f$ definido por:
    \[
    \nabla f(x,y) = \langle f_x(x,y),\, f_y(x,y) \rangle
                  = \frac{\partial f}{\partial x}\,\mathbf{i}
                    + \frac{\partial f}{\partial y}\,\mathbf{j}.
    \]
    \item Si $f(x,y,z)$ es una función de tres variables, su gradiente es:
    \[
    \nabla f(x,y,z) = \langle f_x(x,y,z),\, f_y(x,y,z),\,
    f_z(x,y,z) \rangle
    = \frac{\partial f}{\partial x}\,\mathbf{i}
      + \frac{\partial f}{\partial y}\,\mathbf{j}
      + \frac{\partial f}{\partial z}\,\mathbf{k}.
    \]
\end{itemize}
\end{definition}

\begin{rem}
El símbolo $\nabla$ se lee \emph{nabla} o \emph{del}. El gradiente
transforma una función escalar en un campo vectorial: en cada punto del
dominio, asigna un vector cuyas componentes son las derivadas parciales
de $f$. Este vector contiene toda la información de primer orden sobre
el comportamiento de $f$ en ese punto.
\end{rem}

\begin{example}[Cálculo del gradiente]
Encuentre el gradiente de $f(x,y) = x^2 \sen y + e^{xy}$.
\end{example}

\begin{myproof}
Aplicando la definición, se calculan las derivadas parciales:
\[
f_x(x,y) = 2x \sen y + y e^{xy}, \qquad
f_y(x,y) = x^2 \cos y + x e^{xy}.
\]
Por tanto:
\[
\nabla f(x,y) = \langle 2x \sen y + y e^{xy},\; x^2 \cos y + x e^{xy} \rangle.
\]
En notación de vectores unitarios:
\[
\nabla f(x,y) = (2x \sen y + y e^{xy})\mathbf{i}
              + (x^2 \cos y + x e^{xy})\mathbf{j}. \qquad\square
\]
\end{myproof}

El operador gradiente satisface propiedades algebraicas que lo convierten
en una herramienta flexible. Estas propiedades se deducen directamente de
la linealidad de la derivación parcial y de las reglas del producto y
cociente para funciones escalares.

\begin{theorem}[Propiedades del gradiente]
Sean $f$ y $g$ funciones escalares diferenciables de dos o tres
variables, y sea $k\in\mathbb{R}$ una constante. Entonces:
\begin{enumerate}
    \item \textbf{Linealidad:}
    $\nabla(f \pm g) = \nabla f \pm \nabla g$.
    \item \textbf{Multiplicación por constante:}
    $\nabla(kf) = k\,\nabla f$.
    \item \textbf{Regla del producto:}
    $\nabla(fg) = f\,\nabla g + g\,\nabla f$.
    \item \textbf{Regla del cociente:}
    $\nabla\left(\dfrac{f}{g}\right) =
    \dfrac{g\,\nabla f - f\,\nabla g}{g^2}$, siempre que $g\neq 0$.
\end{enumerate}
\end{theorem}

%%----------------------------------------------------------
\subsection{Interpretación geométrica del gradiente}
%%----------------------------------------------------------

El vector gradiente $\nabla f(x_0,y_0)$ en un punto posee tres
propiedades geométricas fundamentales que lo convierten en el objeto
central del análisis local de funciones de varias variables.

\begin{enumerate}
    \item \textbf{Dirección de mayor aumento:} $\nabla f(x_0,y_0)$
    apunta en la dirección en la que la función $f$ crece más
    rápidamente desde el punto $(x_0,y_0)$.
    \item \textbf{Magnitud:} La longitud del vector gradiente,
    $|\nabla f(x_0,y_0)|$, es la tasa máxima de cambio de $f$ en ese
    punto.
    \item \textbf{Perpendicularidad:} $\nabla f(x_0,y_0)$ es ortogonal
    (perpendicular) a las curvas de nivel de $f$, es decir, a las curvas
    donde $f(x,y) = k$ para alguna constante $k$.
\end{enumerate}

\begin{rem}
La propiedad de perpendicularidad merece atención: si una curva de nivel
está parametrizada por $\mathbf{r}(t) = (x(t), y(t))$ con
$f(\mathbf{r}(t)) = k$, entonces derivando respecto a $t$ se obtiene
$\nabla f(\mathbf{r}(t))\cdot \mathbf{r}'(t) = 0$, lo que prueba que el
gradiente es normal a la curva de nivel en cada punto. Esta observación
será clave más adelante para construir planos tangentes a superficies
dadas en forma implícita.
\end{rem}

%%----------------------------------------------------------
\subsection{Derivada direccional}
%%----------------------------------------------------------

Con el gradiente como herramienta auxiliar, se puede ahora definir la
derivada direccional. Intuitivamente, la derivada direccional mide la
tasa de cambio de $f$ cuando nos movemos desde un punto dado en una
dirección específica, representada por un vector unitario.

\begin{definition}[Derivada direccional]
La \textbf{derivada direccional} de $f(x,y)$ en el punto
$(x_0,y_0)$ y en la dirección del vector unitario
$\mathbf{u} = \langle u_1, u_2 \rangle$ está dada por:
\[
D_{\mathbf{u}} f(x_0,y_0)
= \lim_{h \to 0} \frac{f(x_0 + h u_1,\, y_0 + h u_2) - f(x_0,y_0)}{h},
\]
siempre que el límite exista.
\end{definition}

\begin{rem}
La derivada direccional generaliza las derivadas parciales: si
$\mathbf{u} = \mathbf{i} = \langle 1,0 \rangle$, entonces
$D_{\mathbf{i}} f = f_x$; si $\mathbf{u} = \mathbf{j} = \langle 0,1
\rangle$, entonces $D_{\mathbf{j}} f = f_y$. La exigencia de que
$\mathbf{u}$ sea unitario es crucial: si no lo fuera, el cociente
incremental no mediría tasa de cambio por unidad de longitud recorrida.
\end{rem}

Geométricamente, el vector $\mathbf{u}$ determina una recta $L$ que
contiene al punto $(x_0,y_0)$. El plano que contiene a $L$ y es
perpendicular al plano $xy$ corta a la superficie $z=f(x,y)$ en una
curva $C$. La recta tangente a la curva $C$ en el punto
$(x_0,y_0,f(x_0,y_0))$ tiene como pendiente la derivada direccional
$D_{\mathbf{u}} f(x_0,y_0)$.

El siguiente teorema proporciona la fórmula práctica para calcular la
derivada direccional sin recurrir al límite.

\begin{theorem}[Fórmula del gradiente para la derivada direccional]
Sea $f$ una función diferenciable en el punto $p=(a,b)$, y sea
$\mathbf{u}=\langle u_1,u_2\rangle$ un vector unitario. Entonces:
\[
D_{\mathbf{u}} f(p) = \nabla f(p) \cdot \mathbf{u}
= f_x(a,b) u_1 + f_y(a,b) u_2.
\]
Si el vector unitario forma un ángulo $\theta$ con el eje $x$ positivo,
entonces $\mathbf{u} = \langle \cos\theta, \sen\theta \rangle$ y la
fórmula se convierte en:
\[
D_{\mathbf{u}} f(a,b) = f_x(a,b)\cos\theta + f_y(a,b)\sen\theta.
\]
\end{theorem}

\begin{rem}
La diferenciabilidad de $f$ en $p$ es una hipótesis suficiente pero no
necesaria: la derivada direccional puede existir en todas las direcciones
sin que la función sea diferenciable. Sin embargo, si las derivadas
parciales existen y son continuas en una vecindad de $p$, entonces $f$ es
diferenciable y la fórmula se aplica sin reservas.
\end{rem}

\begin{pasos}[Cálculo de la derivada direccional]
Para calcular la derivada direccional de $f$ en el punto $p$ en la
dirección del vector $\mathbf{v}$ (no necesariamente unitario):
\begin{enumerate}
    \item \textbf{Gradiente.} Calcular el gradiente $\nabla f(x,y)$.
    \item \textbf{Evaluación.} Evaluar el gradiente en el punto $p$:
    $\nabla f(p)$.
    \item \textbf{Normalización.} Normalizar el vector dirección:
    $\mathbf{u} = \dfrac{\mathbf{v}}{|\mathbf{v}|}$.
    \item \textbf{Producto escalar.} Calcular
    $D_{\mathbf{u}} f(p) = \nabla f(p)\cdot \mathbf{u}$.
\end{enumerate}
\end{pasos}

\begin{example}[Derivada direccional]
Encuentre la derivada direccional de $f(x,y) = e^x \sen y$ en el punto
$p = (0, \pi/4)$ en la dirección del vector $\mathbf{v} = \mathbf{i} +
\sqrt{3}\,\mathbf{j}$.
\end{example}

\begin{myproof}
Se siguen los pasos del protocolo de cálculo de la derivada direccional.

\textbf{Paso 1. Gradiente.}
\[
\nabla f(x,y) = \langle e^x \sen y,\; e^x \cos y \rangle.
\]

\textbf{Paso 2. Evaluación en $p=(0,\pi/4)$.}
\[
\nabla f(0,\pi/4)
= \left\langle \sen\left(\frac{\pi}{4}\right),\,
         \cos\left(\frac{\pi}{4}\right) \right\rangle
= \left\langle \frac{\sqrt{2}}{2},\, \frac{\sqrt{2}}{2} \right\rangle.
\]

\textbf{Paso 3. Normalización de $\mathbf{v}$.}
\[
|\mathbf{v}| = \sqrt{1^2 + (\sqrt{3})^2} = \sqrt{4} = 2,
\]
\[
\mathbf{u} = \frac{\mathbf{v}}{|\mathbf{v}|}
= \frac{1}{2}\mathbf{i} + \frac{\sqrt{3}}{2}\mathbf{j}
= \left\langle \frac{1}{2},\, \frac{\sqrt{3}}{2} \right\rangle.
\]

\textbf{Paso 4. Producto escalar.}
\[
D_{\mathbf{u}} f(0,\pi/4)
= \nabla f(0,\pi/4)\cdot \mathbf{u}
= \frac{\sqrt{2}}{2}\cdot\frac{1}{2}
  + \frac{\sqrt{2}}{2}\cdot\frac{\sqrt{3}}{2}
= \frac{\sqrt{2}}{4} + \frac{\sqrt{6}}{4}
= \frac{\sqrt{2}+\sqrt{6}}{4}.
\]
Por tanto:
\[
\boxed{D_{\mathbf{u}} f(0,\pi/4) = \frac{\sqrt{2}+\sqrt{6}}{4}}.
\qquad\square
\]
\end{myproof}

%%----------------------------------------------------------
\subsection{Maximización de la derivada direccional}
%%----------------------------------------------------------

Supongamos que tenemos una función $f$ de dos o tres variables y
queremos saber cómo cambia su valor en distintas direcciones a partir de
un punto. Surgen dos preguntas importantes: ¿en qué dirección cambia más
rápido la función y cuál es esa máxima tasa de cambio? El siguiente
teorema responde ambas preguntas.

\begin{theorem}[Máximo de la derivada direccional]
Supongamos que $f$ es una función diferenciable de dos o tres variables.
El valor máximo de la derivada direccional $D_{\mathbf{u}} f(\mathbf{x})$
es $|\nabla f(\mathbf{x})|$ y se presenta cuando $\mathbf{u}$ tiene la
misma dirección que el vector gradiente $\nabla f(\mathbf{x})$. El valor
mínimo es $-|\nabla f(\mathbf{x})|$ y ocurre en la dirección opuesta
$-\nabla f(\mathbf{x})$.
\end{theorem}

\begin{proof}
De la fórmula $D_{\mathbf{u}} f = \nabla f \cdot \mathbf{u}$ y la
definición de producto escalar:
\[
D_{\mathbf{u}} f = |\nabla f|\,|\mathbf{u}|\cos\theta
= |\nabla f|\cos\theta,
\]
donde $\theta$ es el ángulo entre $\nabla f$ y $\mathbf{u}$. Como
$\mathbf{u}$ es unitario, $|\mathbf{u}|=1$. El máximo de $\cos\theta$ es
$1$, que ocurre cuando $\theta=0$, es decir, cuando $\mathbf{u}$ apunta
en la misma dirección que $\nabla f$. En ese caso
$D_{\mathbf{u}} f = |\nabla f|$. El mínimo es $-1$, que ocurre cuando
$\theta=\pi$, es decir, $\mathbf{u}$ en dirección opuesta al gradiente,
y entonces $D_{\mathbf{u}} f = -|\nabla f|$. $\square$
\end{proof}

\begin{rem}
El teorema tiene una consecuencia práctica importante: si se desea
aumentar el valor de $f$ lo más rápidamente posible, hay que moverse en
la dirección del gradiente; si se desea disminuirlo, en la dirección
opuesta. Esta observación es la base de los métodos de optimización por
descenso del gradiente, ampliamente utilizados en aprendizaje automático
y en la resolución numérica de problemas de optimización.
\end{rem}

\begin{example}[Maximización de la derivada direccional]
Sea $f(x,y) = x e^y$.
\begin{enumerate}[a)]
    \item Determine la razón de cambio de $f$ en el punto $P(2,0)$ en la
    dirección del vector que va de $P$ a $Q\left(-\frac12,\,2\right)$.
    \item ¿En qué dirección $f$ tiene la máxima razón de cambio?
    ¿Cuál es esa máxima razón de cambio?
\end{enumerate}
\end{example}

\begin{myproof}
\textbf{a) Primero se calcula el gradiente:}
\[
\nabla f(x,y) = \langle f_x, f_y \rangle = \langle e^y,\, x e^y \rangle.
\]
Evaluando en $P(2,0)$:
\[
\nabla f(2,0) = \langle e^0,\, 2e^0 \rangle = \langle 1, 2 \rangle.
\]
El vector que va de $P$ a $Q$ es:
\[
\overrightarrow{PQ} = Q - P = \left(-\frac12 - 2,\; 2 - 0\right)
= \left(-\frac{5}{2},\, 2\right).
\]
Normalizando:
\[
\left|\overrightarrow{PQ}\right|
= \sqrt{\left(-\frac{5}{2}\right)^2 + 2^2}
= \sqrt{\frac{25}{4} + 4}
= \sqrt{\frac{41}{4}}
= \frac{\sqrt{41}}{2}.
\]
\[
\mathbf{u} = \frac{\overrightarrow{PQ}}{\left|\overrightarrow{PQ}\right|}
= \frac{2}{\sqrt{41}}\left(-\frac{5}{2},\, 2\right)
= \left(-\frac{5}{\sqrt{41}},\, \frac{4}{\sqrt{41}}\right).
\]
La razón de cambio de $f$ en la dirección de $P$ a $Q$ es:
\[
D_{\mathbf{u}} f(2,0)
= \nabla f(2,0)\cdot \mathbf{u}
= \langle 1,2\rangle \cdot
  \left\langle -\frac{5}{\sqrt{41}},\, \frac{4}{\sqrt{41}} \right\rangle
= -\frac{5}{\sqrt{41}} + \frac{8}{\sqrt{41}}
= \frac{3}{\sqrt{41}}.
\]

\textbf{b) Según el teorema, $f$ se incrementa más rápido en la dirección
del vector gradiente $\nabla f(2,0) = \langle 1,2\rangle$. La razón de
cambio máxima es:}
\[
|\nabla f(2,0)| = |\langle 1,2\rangle| = \sqrt{1^2 + 2^2} = \sqrt{5}.
\]

Por tanto:
\[
\boxed{D_{\mathbf{u}} f(2,0) = \frac{3}{\sqrt{41}},\qquad
\text{dirección de máximo crecimiento: } \langle 1,2\rangle,\qquad
\text{tasa máxima: } \sqrt{5}}.
\qquad\square
\]
\end{myproof}

%%----------------------------------------------------------
\subsection{Plano tangente y recta normal}
%%----------------------------------------------------------

En la lección anterior se introdujo el plano tangente a una superficie
determinada por una función explícita $z=f(x,y)$. Ahora se extiende este
concepto a superficies dadas en forma implícita $F(x,y,z)=k$. La
observación clave es que el gradiente $\nabla F$ es normal a la
superficie en cada punto.

\begin{rem}
Si la superficie está dada por $z=f(x,y)$, se puede expresar en forma
implícita como $F(x,y,z) = f(x,y) - z = 0$. Entonces
$\nabla F = \langle f_x, f_y, -1 \rangle$, que es un vector normal al
plano tangente. La ecuación
\[
f_x(x_0,y_0)(x-x_0) + f_y(x_0,y_0)(y-y_0) - (z-z_0) = 0
\]
recupera la fórmula del plano tangente para funciones explícitas
estudiada en la lección anterior.
\end{rem}

\begin{definition}[Plano tangente y recta normal para superficie implícita]
Sea $F(x,y,z)=k$ una superficie y supóngase que $F$ es diferenciable en
el punto $P(x_0,y_0,z_0)$, con $\nabla F(x_0,y_0,z_0) \neq \mathbf{0}$.
El \textbf{plano tangente} a la superficie en $P$ es el plano que pasa
por $P$ y es perpendicular a $\nabla F(P)$. La \textbf{recta normal} a
la superficie en $P$ es la recta que pasa por $P$ y tiene como vector
director $\nabla F(P)$.
\end{definition}

La ecuación del plano tangente en $(x_0,y_0,z_0)$ es:
\[
\nabla F(x_0,y_0,z_0) \cdot \langle x-x_0,\, y-y_0,\, z-z_0 \rangle = 0.
\]
Equivalentemente:
\[
F_x(x_0,y_0,z_0)(x-x_0)
+ F_y(x_0,y_0,z_0)(y-y_0)
+ F_z(x_0,y_0,z_0)(z-z_0) = 0.
\]

Las ecuaciones paramétricas de la recta normal son:
\[
x = x_0 + t\,F_x(x_0,y_0,z_0),\quad
y = y_0 + t\,F_y(x_0,y_0,z_0),\quad
z = z_0 + t\,F_z(x_0,y_0,z_0),\quad t\in\mathbb{R}.
\]

\begin{pasos}[Plano tangente a superficie implícita]
Para hallar el plano tangente y la recta normal a la superficie
$F(x,y,z)=k$ en el punto $P(x_0,y_0,z_0)$:
\begin{enumerate}
    \item \textbf{Función implícita.} Escribir la superficie en la forma
    $F(x,y,z)=k$ (si es necesario, restando $k$ de ambos lados).
    \item \textbf{Gradiente.} Calcular $\nabla F(x,y,z)$.
    \item \textbf{Evaluación.} Evaluar $\nabla F$ en $P$.
    \item \textbf{Ecuaciones.} Escribir la ecuación del plano usando
    $\nabla F(P)$ como vector normal, y las ecuaciones paramétricas de la
    recta usando $\nabla F(P)$ como vector director.
\end{enumerate}
\end{pasos}

\begin{example}[Plano tangente y recta normal a superficie implícita]
Halle las ecuaciones del plano tangente y la recta normal a la superficie
$x e^{yz} = 1$ en el punto $(1,0,5)$.
\end{example}

\begin{myproof}
Se sigue el protocolo para el plano tangente a superficie implícita.

\textbf{Paso 1. Función implícita.} Se define
\[
F(x,y,z) = x e^{yz} - 1 = 0.
\]
La superficie es $F(x,y,z)=0$.

\textbf{Paso 2. Gradiente.}
\[
\nabla F(x,y,z)
= \langle F_x, F_y, F_z \rangle
= \langle e^{yz},\, xz e^{yz},\, xy e^{yz} \rangle.
\]

\textbf{Paso 3. Evaluación en $P(1,0,5)$.}
\[
\nabla F(1,0,5)
= \langle e^{0\cdot 5},\, 1\cdot 5\cdot e^{0\cdot 5},\,
  1\cdot 0\cdot e^{0\cdot 5} \rangle
= \langle 1,\, 5,\, 0 \rangle.
\]

\textbf{Paso 4. Ecuaciones.}
El plano tangente tiene vector normal $\mathbf{n} = \langle 1,5,0 \rangle$
y pasa por $(1,0,5)$:
\[
1(x-1) + 5(y-0) + 0(z-5) = 0
\implies x - 1 + 5y = 0
\implies x + 5y = 1.
\]
La recta normal tiene vector director $\mathbf{v} = \langle 1,5,0 \rangle$
y pasa por $(1,0,5)$:
\[
x = 1 + t,\qquad y = 5t,\qquad z = 5,\qquad t\in\mathbb{R}.
\]
Por tanto:
\[
\boxed{\text{Plano tangente: } x + 5y = 1,\qquad
\text{Recta normal: } (x,y,z) = (1,0,5) + t\langle 1,5,0 \rangle}.
\qquad\square
\]
\end{myproof}

%%==========================================================
\section{Valores extremos de una función multivariable}
%%==========================================================

La teoría de máximos y mínimos de funciones de una variable se extiende
naturalmente a funciones de varias variables, aunque la geometría es más
rica: los puntos críticos pueden ser máximos, mínimos o puntos de silla,
una configuración sin análogo en una dimensión.

%%----------------------------------------------------------
\subsection{Definiciones básicas}
%%----------------------------------------------------------

\begin{definition}[Extremos absolutos y locales]
Sea $f$ una función de dos variables con dominio $S$, y sea
$p_0\in S$.
\begin{itemize}
    \item $f(p_0)$ es un \textbf{valor máximo absoluto} de $f$ en $S$ si
    $f(p_0) \geq f(p)$ para todo $p\in S$.
    \item $f(p_0)$ es un \textbf{valor mínimo absoluto} de $f$ en $S$ si
    $f(p_0) \leq f(p)$ para todo $p\in S$.
    \item $f(p_0)$ es un \textbf{valor máximo local} (o relativo) si
    $f(p_0) \geq f(p)$ para todo $p$ suficientemente cercano a $p_0$.
    \item $f(p_0)$ es un \textbf{valor mínimo local} (o relativo) si
    $f(p_0) \leq f(p)$ para todo $p$ suficientemente cercano a $p_0$.
\end{itemize}
\end{definition}

\begin{theorem}[Condición necesaria para extremo local]
Si $f$ tiene un máximo local o un mínimo local en $(a,b)$ y las
derivadas parciales de primer orden de $f$ existen ahí, entonces
\[
f_x(a,b) = 0 \qquad\text{y}\qquad f_y(a,b) = 0.
\]
\end{theorem}

\begin{proof}
Si $f$ tiene un extremo local en $(a,b)$, entonces la función de una
variable $g(x) = f(x,b)$ tiene un extremo local en $x=a$. Por el
teorema de Fermat para funciones de una variable, $g'(a)=f_x(a,b)=0$.
Análogamente, $h(y)=f(a,y)$ tiene un extremo local en $y=b$, luego
$h'(b)=f_y(a,b)=0$. $\square$
\end{proof}

\begin{definition}[Punto crítico]
Un punto $(a,b)$ en el interior del dominio de $f$ es un
\textbf{punto crítico} si $f_x(a,b)=0$ y $f_y(a,b)=0$, o si al menos
una de las derivadas parciales no existe.
\end{definition}

\begin{rem}
La condición $f_x=f_y=0$ es necesaria pero no suficiente para la
existencia de un extremo local. Un punto donde ambas derivadas parciales
se anulan puede ser un máximo, un mínimo, o un \textbf{punto de silla},
que es un punto donde la función crece en una dirección y decrece en
otra.
\end{rem}

%%----------------------------------------------------------
\subsection{Teorema del Valor Extremo}
%%----------------------------------------------------------

El siguiente resultado garantiza la existencia de extremos absolutos para
funciones continuas sobre conjuntos cerrados y acotados.

\begin{theorem}[Teorema del Valor Extremo de Weierstrass]
Si $f$ es una función continua de dos variables en un conjunto $S$ que es
\textbf{cerrado} y \textbf{acotado}, entonces $f$ alcanza un valor
máximo absoluto y un valor mínimo absoluto en $S$.
\end{theorem}

\begin{rem}
Las hipótesis de cerradura y acotación son esenciales. Si el conjunto no
es cerrado, el extremo puede ocurrir en la frontera que no pertenece al
dominio. Si no es acotado, la función puede crecer o decrecer
indefinidamente. En una variable, el intervalo $[a,b]$ es el análogo de
un conjunto cerrado y acotado en $\mathbb{R}^2$.
\end{rem}

\begin{pasos}[Localización de extremos absolutos en un conjunto cerrado y acotado]
Para encontrar los extremos absolutos de $f$ en un conjunto $S$ cerrado y
acotado:
\begin{enumerate}
    \item \textbf{Puntos críticos interiores.} Encontrar todos los
    puntos críticos de $f$ en el interior de $S$ (donde $f_x=f_y=0$ o no
    existen).
    \item \textbf{Frontera.} Parametrizar la frontera de $S$ y reducir el
    problema a funciones de una variable sobre cada segmento de la
    frontera. Encontrar los puntos críticos de estas funciones
    unidimensionales, junto con los vértices o esquinas.
    \item \textbf{Evaluación.} Evaluar $f$ en todos los puntos candidatos
    obtenidos en los pasos 1 y 2.
    \item \textbf{Comparación.} El mayor valor es el máximo absoluto; el
    menor, el mínimo absoluto.
\end{enumerate}
\end{pasos}

\begin{rem}
En la práctica, el paso 2 (frontera) es el que requiere más atención,
pues la frontera de una región en $\mathbb{R}^2$ suele estar formada por
curvas (segmentos de recta, arcos de circunferencia, etc.) que deben
parametrizarse adecuadamente. Cada segmento de frontera produce una
función de una variable $f(\mathbf{r}(t))$ cuyo dominio es un intervalo
cerrado, donde se aplica el método de localización de extremos en una
variable.
\end{rem}

%%----------------------------------------------------------
\subsection{Criterio de la segunda derivada}
%%----------------------------------------------------------

El criterio de la segunda derivada para funciones de dos variables
clasifica los puntos críticos donde $f_x=f_y=0$ utilizando las derivadas
parciales de segundo orden.

\begin{definition}[Discriminante o Hessiano]
Para una función $f$ de dos variables con derivadas parciales de segundo
orden continuas, se define el \textbf{discriminante} (o Hessiano) en un
punto $(a,b)$ como:
\[
D(a,b) = f_{xx}(a,b) f_{yy}(a,b) - \bigl[f_{xy}(a,b)\bigr]^2.
\]
\end{definition}

\begin{theorem}[Criterio de la segunda derivada]
Sea $f$ una función con derivadas parciales de segundo orden continuas en
una vecindad de un punto crítico $(a,b)$ donde $f_x(a,b)=0$ y
$f_y(a,b)=0$. Sea $D = f_{xx}f_{yy} - (f_{xy})^2$ evaluado en $(a,b)$.
Entonces:
\begin{enumerate}[a)]
    \item Si $D > 0$ y $f_{xx}(a,b) > 0$, entonces $f(a,b)$ es un
    \textbf{mínimo local}.
    \item Si $D > 0$ y $f_{xx}(a,b) < 0$, entonces $f(a,b)$ es un
    \textbf{máximo local}.
    \item Si $D < 0$, entonces $(a,b)$ es un \textbf{punto de silla}.
    \item Si $D = 0$, el criterio no permite concluir nada (el punto
    puede ser máximo, mínimo o punto de silla).
\end{enumerate}
\end{theorem}

\begin{rem}
El discriminante puede recordarse como el determinante de la matriz
hessiana:
\[
D = \begin{vmatrix}
f_{xx} & f_{xy} \\
f_{yx} & f_{yy}
\end{vmatrix}
= f_{xx} f_{yy} - (f_{xy})^2.
\]
El caso $D=0$ es el único que requiere un análisis adicional, ya que la
información de segundo orden no es suficiente para clasificar el punto
crítico.
\end{rem}

\begin{rem}[Seguimiento de hipótesis]
El criterio de la segunda derivada exige que las derivadas parciales de
segundo orden sean continuas en una vecindad del punto crítico. Si esta
condición falla, el criterio no es aplicable. Por ejemplo, la función
$f(x,y) = x^2 y^2$ tiene un punto crítico en $(0,0)$ donde $D=0$, y la
continuidad de las derivadas segundas no es el problema (son continuas),
sino que $D=0$ es un caso inconcluso. En cambio, funciones con
discontinuidades en las derivadas segundas pueden presentar
comportamientos atípicos que el criterio no captura.
\end{rem}

\begin{pasos}[Clasificación de puntos críticos mediante el criterio de la segunda derivada]
Para clasificar los puntos críticos de $f$:
\begin{enumerate}
    \item \textbf{Puntos críticos.} Resolver el sistema
    $f_x=0$, $f_y=0$ para obtener todos los puntos críticos.
    \item \textbf{Derivadas segundas.} Calcular $f_{xx}$, $f_{xy}$,
    $f_{yy}$.
    \item \textbf{Discriminante.} Para cada punto crítico $(a,b)$,
    calcular $D = f_{xx} f_{yy} - (f_{xy})^2$.
    \item \textbf{Clasificación.} Aplicar el criterio de la segunda
    derivada:
    \begin{itemize}
        \item $D>0$ y $f_{xx}>0$: mínimo local.
        \item $D>0$ y $f_{xx}<0$: máximo local.
        \item $D<0$: punto de silla.
        \item $D=0$: criterio inconcluso.
    \end{itemize}
\end{enumerate}
\end{pasos}

\begin{example}[Clasificación de puntos críticos con punto de silla]
Determine los valores máximos y mínimos locales y los puntos de silla de
$f(x,y) = x^4 + y^4 - 4xy + 1$.
\end{example}

\begin{myproof}
Se sigue el protocolo de clasificación de puntos críticos.

\textbf{Paso 1. Puntos críticos.}
\[
f_x = 4x^3 - 4y = 0 \implies y = x^3.
\]
\[
f_y = 4y^3 - 4x = 0 \implies x = y^3.
\]
Sustituyendo $y=x^3$ en $x=y^3$:
\[
x = (x^3)^3 = x^9 \implies x^9 - x = 0 \implies x(x^8 - 1)=0.
\]
Las raíces reales son $x=0$, $x=1$, $x=-1$. Por tanto, los puntos
críticos son:
\[
(0,0),\qquad (1,1),\qquad (-1,-1).
\]

\textbf{Paso 2. Derivadas segundas.}
\[
f_{xx} = 12x^2,\qquad f_{xy} = -4,\qquad f_{yy} = 12y^2.
\]

\textbf{Paso 3. Discriminante.}
\[
D(x,y) = f_{xx} f_{yy} - (f_{xy})^2
= (12x^2)(12y^2) - (-4)^2
= 144x^2 y^2 - 16.
\]

\textbf{Paso 4. Clasificación.}

\textbf{En $(0,0)$:}
\[
D(0,0) = 0 - 16 = -16 < 0.
\]
Por el caso c) del criterio, $(0,0)$ es un \textbf{punto de silla}.
El valor en el punto es $f(0,0)=1$.

\textbf{En $(1,1)$:}
\[
D(1,1) = 144(1)(1) - 16 = 128 > 0,
\]
\[
f_{xx}(1,1) = 12 > 0.
\]
Por el caso a), $f(1,1) = 1 + 1 - 4 + 1 = -1$ es un
\textbf{mínimo local}.

\textbf{En $(-1,-1)$:}
\[
D(-1,-1) = 144(1)(1) - 16 = 128 > 0,
\]
\[
f_{xx}(-1,-1) = 12 > 0.
\]
Por el caso a), $f(-1,-1) = 1 + 1 - 4 + 1 = -1$ es un
\textbf{mínimo local}.

Por tanto:
\[
\boxed{
\text{Punto de silla: } (0,0,1),\quad
\text{Mínimos locales: } (1,1,-1),\ (-1,-1,-1)
}.
\qquad\square
\]
\end{myproof}

\begin{rem}
El ejemplo muestra un fenómeno que no ocurre en funciones de una
variable: el punto de silla. En $(0,0)$, la función crece en las
direcciones de los ejes (pues $x^4$ y $y^4$ son positivos) pero decrece
en las direcciones diagonales (pues el término $-4xy$ domina cerca del
origen). Esta geometría se asemeja a una silla de montar: el punto
central es un mínimo en una dirección y un máximo en otra, sin ser ni uno
ni otro.
\end{rem}

%%==========================================================
\section{Problemas Propuestos}
%%==========================================================

\begin{prob}[Gradiente y derivada direccional]
Calcule el gradiente de las siguientes funciones:
\begin{multicols}{2}
\begin{enumerate}
    \item $f(x,y) = x^2 y + \sen(xy)$
    \item $f(x,y) = \ln(x^2 + y^2)$
    \item $f(x,y,z) = xyz + e^{xz}$
    \item $f(x,y,z) = \arctan(x/y) + z^2$
\end{enumerate}
\end{multicols}
\end{prob}

\begin{prob}
Encuentre la derivada direccional de $f$ en el punto $p$ en la dirección
del vector $\mathbf{v}$:
\begin{enumerate}[a)]
    \item $f(x,y) = x^2 + y^2$, $p=(1,2)$, $\mathbf{v} = \langle 3,4 \rangle$
    \item $f(x,y) = e^x \cos y$, $p=(0,\pi/3)$, $\mathbf{v} = \langle 1,-1 \rangle$
    \item $f(x,y) = \sqrt{xy}$, $p=(4,9)$, $\mathbf{v} = \langle -1,2 \rangle$
\end{enumerate}
\end{prob}

\begin{prob}
Para cada función, determine la dirección de máximo crecimiento y la tasa
máxima de cambio en el punto dado:
\begin{enumerate}[a)]
    \item $f(x,y) = 3x^2 - 2xy + y^2$, $p=(1,2)$
    \item $f(x,y) = \sinh(xy)$, $p=(\pi/2,1)$
    \item $f(x,y,z) = x^2 + y^2 + z^2 - xy$, $p=(1,1,1)$
\end{enumerate}
\end{prob}

\begin{prob}
Encuentre la ecuación del plano tangente y la recta normal a la
superficie en el punto indicado:
\begin{enumerate}[a)]
    \item $x^2 + y^2 + z^2 = 14$, $(1,2,3)$
    \item $x^2 + 2y^2 - z^2 = 4$, $(2,1,-1)$
    \item $xy + yz + zx = 1$, $(1,0,1)$
\end{enumerate}
\end{prob}

\begin{prob}[Extremos absolutos]
Encuentre los extremos absolutos de $f$ en el conjunto $S$ indicado:
\begin{enumerate}[a)]
    \item $f(x,y) = x^2 + y^2 - 2x - 4y$, $S$: el rectángulo
    $0\leq x\leq 3$, $0\leq y\leq 2$
    \item $f(x,y) = xy$, $S$: el círculo $x^2 + y^2 \leq 1$
    \item $f(x,y) = x^2 + 2y^2$, $S$: el triángulo con vértices
    $(0,0)$, $(1,0)$, $(0,1)$
\end{enumerate}
\end{prob}

\begin{prob}[Clasificación de puntos críticos]
Encuentre y clasifique todos los puntos críticos de las siguientes
funciones:
\begin{multicols}{2}
\begin{enumerate}
    \item $f(x,y) = x^2 + xy + y^2 - 3x$
    \item $f(x,y) = 3x^2 y + y^3 - 3y$
    \item $f(x,y) = x^3 + y^3 - 3xy$
    \item $f(x,y) = x^2 y^2 - 8x^2 - 4y^2 + 24$
    \item $f(x,y) = x^4 + y^4 - 4x^2 y^2$
    \item $f(x,y) = x^3 + 3xy^2 - 15x - 12y$
\end{enumerate}
\end{multicols}
\end{prob}

\begin{prob}
Para la función $f(x,y) = x^3 - 3xy^2$:
\begin{enumerate}[a)]
    \item Encuentre todos los puntos críticos.
    \item Clasifique cada punto crítico usando el criterio de la segunda
    derivada.
    \item Verifique que el origen es un punto de silla mostrando que
    $f$ crece en una dirección y decrece en otra.
\end{enumerate}
\end{prob}

\begin{prob}
Demuestre que la función $f(x,y) = x^2 y^2$ tiene un punto crítico en
$(0,0)$ donde el criterio de la segunda derivada es inconcluso.
Determine la naturaleza del punto mediante un análisis directo de los
valores de $f$ en una vecindad del origen.
\end{prob}

\begin{prob}
Una función $f$ tiene derivadas parciales continuas hasta segundo orden.
En el punto crítico $(1,2)$ se sabe que $f_{xx}=4$, $f_{yy}=9$ y
$f_{xy}=6$. ¿Qué puede concluirse sobre la naturaleza del punto
crítico?
\end{prob}

\begin{prob}
Encuentre los puntos de la superficie $z = x^2 + y^2$ donde el plano
tangente es paralelo al plano $2x + 4y - z = 0$.
\end{prob}

% ============================================================
% Cap. 20 — Diferenciación parcial, gradiente y diferenciabilidad
% 9 problemas unicos: 9 listos para integrar ahora, 0 esperan pasada Claude Code (pstricks)
% ============================================================

% >>> LISTOS PARA INTEGRAR AHORA <<<

% [Original: líneas 459-478 | solución completa | sin figura]
\begin{prob}
Demuestre que el elipsoide $3x^2+2y^2+z^2=9$ y la esfera $x^2+y^2+z^2-8x-6y-8z+24=0$ son tangentes entre sí en el punto $(1,1,2)$. (Esto quiere decir que tienen un plano tangente común en ese punto.)
\begin{myproof}
Sea $F(x,y,z)=3x^2+2y^2+z^2-9,$ entonces $$\nabla F=6x{\bf i}+4y{\bf j}+2z{\bf k}$$ es su vector gradiente. Como el plano tangente es igual a $$\nabla F \cdot \left\langle x-x_0,y-y_0,z-z_0\right\rangle=0,$$ entonces el plano tangente tiene ecuación
$$ 0=6x_0(x-x_0)+4y_0(y-y_0)+2z_0(z-z_0).$$  
Si $(x_0,y_0,z_0)=(1,1,2)$ la ecuación del plano tangente al elipsoide en este punto es 
\begin{align*}
0&=6(1)(x-1)+4(y-1)+4(z-2)\\
0&=6(x-1)+4(y-1)+4(z-2).
\end{align*}
Análogamente, sea $G(x,y,z)=x^2+y^2+z^2-8x-6y-8z+24,$ entonces $$\nabla G=(2x-8){\bf i}+(2y-6){\bf j}+(2z-8){\bf k}$$ es su vector gradiente. Como el plano tangente es igual a $$\nabla G \cdot \left\langle x-x_0,y-y_0,z-z_0\right\rangle=0,$$ entonces el plano tangente tiene ecuación
$$ 0=(2x_0-8)(x-x_0)+(2y_0-6)(y-y_0)+(2z-8)(z-z_0).$$  
Si $(x_0,y_0,z_0)=(1,1,2)$ la ecuación del plano tangente a la esfera en este punto es 
\begin{align*}
0&=(2(1)-8)(x-1)+(2(1)-6)(y-1)+(2(2)-8)(z-2)\\
0&=-6(x-1)-4(y-1)-4(z-2).
\end{align*}
Finalmente observe que los vectores normales a cada plano tangente son paralelos, por lo tanto, son el mismo plano tangente y así se cumple que las superficies son tangentes entre sí en el punto $(1,1,2).$
\end{myproof}
\end{prob}

% [Original: líneas 715-723 | solución completa | sin figura]
\begin{prob}
Determine el punto del paraboloide $y=x^2+z^2$ donde el plano tangente es paralelo al plano $x+2y+3z=1.$
\begin{myproof}
Sea $F(x,y,z)=x^2+z^2-y,$ entonces $y=x^2+z^2$ es una superficie de nivel de $F.$ Note que $\nabla F=\left\langle 2x,-1,2z \right\rangle$ es normal a la superficie, por lo tanto, es normal al plano tangente. Como el plano tangente es paralelo al plano $x+2y+3z=1$ entonces los vectores normales de cada plano deben ser paralelos. Así, $\left\langle 2x,-1,2z \right\rangle=k\left\langle 1,2,3 \right\rangle.$
Despejando, $2x=k,$ $-1=2k$ y $2z=3k,$ lo cual implica que $k=\dfrac{-1}{2}.$ 
Si $(x_0,y_0,z_0)$ es un punto en el paraboloide debe cumplir que $$\left\langle 2x_0,-1,2z_0 \right\rangle=\left\langle -1/2,-1,-3/2 \right\rangle.$$
Finalmente, si $x_0=-1/4$ y $z_0=-3/4$ al despejar $y_0$ de la ecuación queda $$y_0=\left( \frac{-1}{4} \right)^2+\left( \frac{-3}{4} \right)^2=\frac{5}{8}.$$ El punto es $\left( \dfrac{-1}{4},\dfrac{5}{8},\dfrac{-3}{4} \right).$
\end{myproof}
\end{prob}

% [Original: líneas 876-882 | solución completa | sin figura]
\begin{prob}
Halle los puntos del hiperboloide $x^2-y^2-z^2=1$ donde el plano tangente es paralelo al plano $z = x + y.$
\begin{myproof}
Sea $F(x,y,z)=x^2-y^2-z^2-1,$ entonces $$\nabla F=2x{\bf i}-2y{\bf j}-2z{\bf k}$$ es su vector gradiente. Como el plano tangente es paralelo al plano $x+y-z=0,$ sus vectores normales deben ser paralelos y así $ \left\langle 2x_0,-2y_0,-2z_0\right\rangle =k \left\langle 1,1,-1\right\rangle. $ Lo cual es equivalente a $$ \left\langle x_0,-y_0,-z_0\right\rangle =k \left\langle 1,1,-1\right\rangle. $$
Esto implica que  $x_0=k$, $y_0=-k$ y $z_0=k.$ Reemplazando en la ecuación del hiperboloide $k^2-(-k)^2-k^2=1,$ lo cual implica que $k^2=-1$ lo cual es imposible. Así, no existe tal punto sobre el hiperboloide.
\end{myproof}
\end{prob}

% [Original: líneas 929-931 | sin verificar | sin figura | 2 variante(s) de parámetro (se descartaron 1)]
\begin{prob}
Encuentre la ecuación del plano tangente a la superficie $xy+yz+zx=7$ en el punto $(-1,3,5).$
\end{prob}

% [Original: líneas 949-951 | sin verificar | sin figura]
\begin{prob}
Encuentre la ecuación del plano tangente a la superficie $x^2+2y^2+4z^2=9$ en el punto $(2,1,\sqrt{3}/2).$
\end{prob}

% [Original: líneas 989-991 | sin verificar | sin figura | 2 variante(s) de parámetro (se descartaron 1)]
\begin{prob}
Encuentre la ecuación del plano tangente a la superficie $5x^2-y^2+4z^2=8$ en el punto $(-2,4,1).$
\end{prob}

% [Original: líneas 1224-1226 | sin verificar | sin figura | 3 variante(s) de parámetro (se descartaron 2)]
\begin{prob}
Considere la función $f(x,y)=3x^2y^4-12x^6+2xy^5.$ Calcule $\dfrac{\partial f}{\partial x}$ y $\dfrac{\partial f}{\partial y}$ para verificar que \[  x\left( \dfrac{\partial f}{\partial x} \right)+y\left( \dfrac{\partial f}{\partial y} \right)=6f(x,y).\] 
\end{prob}

% [Original: líneas 1252-1254 | sin verificar | sin figura]
\begin{prob}
Una función de dos variables $f(x,y)$ que satisface la ecuación de Laplace, \[\dfrac{\partial^2f}{\partial x^2}+\dfrac{\partial^2f}{\partial y^2}=0\] se dice que es armónica. Determine si la función $f(x,y)=Ay^3+5x$ es armónica.
\end{prob}

% [Original: líneas 1408-1410 | sin verificar | sin figura]
\begin{prob}
Encuentre la ecuación del plano tangente a la superficie $x^2+2y^2+z^2=4$ en el punto $(1,-1,1).$
\end{prob}

% ============================================================
% Cap. 21 — Regla de la cadena y derivadas direccionales
% 17 problemas unicos: 15 listos para integrar ahora, 2 esperan pasada Claude Code (pstricks)
% ============================================================

% >>> LISTOS PARA INTEGRAR AHORA <<<

% [Original: líneas 448-457 | solución completa | sin figura]
\begin{prob}
El voltaje $V$ en un circuito eléctrico simple disminuye con lentitud a medida que la batería se gasta. La resistencia $R$ se incrementa lentamente cuando el resistor se calienta. Mediante la ley de Ohm, $V=IR$, determine cómo cambia la corriente $I$ en el momento en que $R=500\ \Omega$ , $I=0.06\  A$, $\frac{dV}{dt}=-0.02\ V/s$ y $\frac{dR}{dt}=0.04\ \Omega /s.$
\begin{myproof}
El voltaje es $V=IR.$ Se conoce que $\dfrac{dV}{dt}=-0.02 \ V/s$ y $\dfrac{dR}{dt}=0.04\ \Omega/s.$ Si $R=500\ \Omega$ e $I=0.06\ A$ al usar regla de la cadena se obtiene que,
\begin{align*}
\dfrac{dV}{dt}&=\dfrac{\partial V}{\partial I}\dfrac{dI}{dt}+\dfrac{\partial V}{\partial R}\dfrac{dR}{dt}\\&=\left(R\right)\left(\dfrac{dI}{dt}\right)+\left(I\right)\left(\dfrac{dR}{dt}\right),\\
-0.02&=\left(500\right)\left(\dfrac{dI}{dt}\right)+\left(0.06\right)\left(0.04\right),\\-0.02&=\left( 500\dfrac{dI}{dt} \right)+\left( \dfrac{3}{1250}  \right)\\-0.02-\left( \dfrac{3}{1250}  \right)&=\left( 500\dfrac{dI}{dt} \right)\\-\dfrac{14}{625}&=\left( 500\dfrac{dI}{dt} \right)\\-\dfrac{14}{312500}\ A/s&=\dfrac{dI}{dt}\\\dfrac{dI}{dt}&\approx -0.0000448\ A/s.
\end{align*}
\end{myproof}
\end{prob}

% [Original: líneas 704-713 | solución completa | sin figura]
\begin{prob}
El radio de un cono circular recto se incrementa a una razón de $5\ m/s$, mientras su altura disminuye a razón de $10\ m/s$. ¿A qué razón cambia el volumen del cono cuando el radio es $800\ m$ y la altura es de $200\ m$? Construya una representación de la situación.
\begin{myproof}
El volumen del cono es $V=\dfrac{\pi}{3}r^2h.$ Se conoce que $\dfrac{dr}{dt}=5 \ m/s$ y $\dfrac{dh}{dt}=-10\ m/s.$ Si $h=200\ m$ y $r=800\ m$ al usar regla de la cadena se obtiene que,
\begin{align*}
\dfrac{dV}{dt}=\dfrac{\partial V}{\partial r}\dfrac{dr}{dt}+\dfrac{\partial V}{\partial h}\dfrac{dh}{dt}&=\left(\dfrac{2}{3}\pi r h\right)\left(5\right)+\left(\dfrac{1}{3}\pi r^2\right)\left(-10\right),\\
&=\left(\dfrac{2}{3}\pi (800) (200)\right)\left(5\right)+\left(\dfrac{1}{3}\pi (800)^2\right)\left(-10\right),\\&=\left(\dfrac{1600000}{3}\pi\right)+\left(\dfrac{-6400000}{3}\pi\right)\\&=\left(\dfrac{-4800000}{3}\pi\right)\ m/s.
\end{align*}
\end{myproof}
\end{prob}

% [Original: líneas 865-874 | solución completa | sin figura]
\begin{prob}
La longitud $l$, ancho $a$ y altura $h$ de una caja cambia con el tiempo. En un cierto instante, las dimensiones son $l= 5\ m,$ $a=3\ m$ y $h= 9\ m$, y se conoce que $a$ y $h$ se incrementan a razón de $0.3\ m/s$, en tanto que $l$ disminuye a razón de $0.5\ m/s.$ Encuentre en ese instante las razones a las cuales el área superficial de la caja cambia.
\begin{myproof}
El área superficial de la caja es $S=2la+2ah+2hl.$ Se conoce que $\dfrac{da}{dt}=0.3=\dfrac{dh}{dt}$ y $\dfrac{dl}{dt}=-0.5\ m/s.$ Si $l=5\ m,$ $a=3\ m,$ y $h=9\ m$ al usar regla de la cadena se obtiene que,
\begin{align*}
\dfrac{dS}{dt}&=\dfrac{\partial S}{\partial l}\dfrac{dl}{dt}+\dfrac{\partial S}{\partial h}\dfrac{dh}{dt}+\dfrac{\partial S}{\partial a}\dfrac{da}{dt}\\&=\left( 2a+2h \right)\left(\dfrac{dl}{dt}\right)+\left(2a+2l\right)\left(\dfrac{dh}{dt}\right)+\left(2h+2l\right)\left(\dfrac{da}{dt}\right),\\
\dfrac{dS}{dt}&=\left(24\right)\left(-0.5\right)+\left(16\right)\left(0.03\right)+\left(28\right)\left(0.03\right),\\\dfrac{dS}{dt}&=6/5\approx 1.2\ m/s.
\end{align*}
\end{myproof}
\end{prob}

% [Original: líneas 925-927 | sin verificar | sin figura | 2 variante(s) de parámetro (se descartaron 1)]
\begin{prob}
La longitud $l$, ancho $a$ y altura $h$ de una caja cambia con el tiempo. En un cierto instante, las dimensiones son $l= 15\ m,$ $a=12\ m$ y $h= 3\ m$, y se conoce que $a$ y $h$ se incrementan a razón de $0.2\ m/s$, en tanto que $l$ disminuye a razón de $0.1\ m/s.$ Encuentre en ese instante las razones a las cuales el volumen de la caja cambia.
\end{prob}

% [Original: líneas 1127-1129 | sin verificar | sin figura | 2 variante(s) de parámetro (se descartaron 1)]
\begin{prob}
La longitud $x,$ el ancho $y$ y la altura $z$ de una caja cambian con el tiempo. En cierto instante las dimensiones de la caja son $x=A\ cm,$ $y=2A\ cm$ y $z= 2A\ cm.$ Si $x$ y $y$ aumentan a una tasa de $2A\ cm/s$ mientras que $z$ disminuye a razón de $3A\ m/s,$ emplee la regla de la cadena para determinar la tasa a la cual cambia el volumen de la caja en ese instante.
\end{prob}

% [Original: líneas 1131-1133 | sin verificar | sin figura | 2 variante(s) de parámetro (se descartaron 1)]
\begin{prob}
La temperatura de una colina en un punto $(x,y,z)$ está dada por la función $$T(x,y,z)=200e^{-x^2-3y^2-9z^2}$$ donde $T$ se mide en grados Celsius y $x,\ y, \ z$ en metros.  Determine la dirección que un insecto seguiría, empezando en el punto $(2A,-A,2A)$ con el fin de calentarse lo más rápido posible. Justifique su respuesta.
\end{prob}

% [Original: líneas 1139-1141 | sin verificar | sin figura | 2 variante(s) de parámetro (se descartaron 1)]
\begin{prob}
La longitud $x,$ el ancho $y$ y la altura $z$ de una caja cambian con el tiempo. En cierto instante las dimensiones de la caja son $x=7A\ cm,$ $y=13A\ cm$ y $z= 11A\ cm.$ Si $x$ y $y$ aumentan a una tasa de $2A\ cm/s$ mientras que $z$ disminuye a razón de $5A\ m/s,$ emplee la regla de la cadena para determinar la tasa a la cual cambia el área superficial de la caja en ese instante (asuma que la caja tiene tapa).
\end{prob}

% [Original: líneas 1143-1145 | sin verificar | sin figura]
\begin{prob}
La temperatura de una colina en un punto $(x,y,z)$ está dada por la función $$T(x,y,z)=236e^{-13x^2-11y^2-17z^2}$$ donde $T$ se mide en grados Celsius y $x,\ y, \ z$ en metros.  Determine la dirección que un insecto seguiría, empezando en el punto $(3A,-7A,11A)$ con el fin de calentarse lo más rápido posible. Justifique su respuesta.
\end{prob}

% [Original: líneas 1156-1158 | sin verificar | sin figura]
\begin{prob}
En un instante dado, la longitud de un cateto de un triángulo rectángulo es de $10\ cm$ y crece a tasa de $1\ cm/min,$ y la longitud del otro cateto es de $12\ cm/min.$ Calcule la tasa de variación de la medida del ángulo opuesto al cateto de 10 cm en ese instante.
\end{prob}

% [Original: líneas 1160-1162 | sin verificar | sin figura]
\begin{prob}
Si $h(x,y)=2e^{-x^2}+3e^{-3y^2}$ representa la altura de una montaña, ¿en qué dirección desde $(1,0,2e^{-1}+3)$ se debería comenzar a caminar para escalar más rápidamente?
\end{prob}

% [Original: líneas 1188-1190 | sin verificar | sin figura]
\begin{prob}
La temperatura de una colina en un punto $(x,y,z)$ está dada por la función $$T(x,y,z)=258e^{-2x^2-5y^2-7z^2}$$ donde $T$ se mide en grados Celsius y $x,\ y, \ z$ en metros.  Determine la dirección que un insecto seguiría, empezando en el punto $(9A,-5A,13A)$ con el fin de calentarse lo más rápido posible. Justifique su respuesta.
\end{prob}

% [Original: líneas 1216-1218 | sin verificar | sin figura]
\begin{prob}
Considere una placa rectangular comprendida por el plano $xy$ cuya temperatura está determinada por  la función $T(x,y,z)=10+8x^2+25y^2$ donde $T$ se mide en grados Celsius y $x,\ y$ en metros.  Determine la dirección que un insecto seguiría, empezando en el punto $(A+1,A)$ con el fin de enfriarse lo más rápido posible. Justifique su respuesta.
\end{prob}

% [Original: líneas 1228-1230 | sin verificar | sin figura | 2 variante(s) de parámetro (se descartaron 1)]
\begin{prob}
La temperatura de una placa en un punto $(x,y)$ está dada por la función $$T(x,y)=10+5x^2-3y^2$$ donde $T$ se mide en grados Celsius y $x,\ y$ en metros.  Determine la dirección que un insecto seguiría, empezando en el punto $(3,-7)$ con el fin de calentarse lo más rápido posible. Justifique su respuesta.
\end{prob}

% [Original: líneas 1260-1262 | sin verificar | sin figura]
\begin{prob}
La temperatura de una placa en un punto $(x,y)$ está dada por la función $$T(x,y)=40+8x^2-25y^2$$ donde $T$ se mide en grados Celsius y $x,\ y$ en metros.  Determine la dirección que un insecto seguiría, empezando en el punto $(9,25)$ con el fin de calentarse lo más rápido posible. Justifique su respuesta.
\end{prob}

% [Original: líneas 1292-1294 | sin verificar | sin figura]
\begin{prob}
Considere una placa rectangular comprendida por el plano $xy$ cuya temperatura está dada por $T(x,y)=5+x^{2}+5y^{2}$ para cada punto $(x,y).$ Determine la dirección que un insecto seguiría, empezando en el punto $(A,A+1)$ con el fin de enfriarse lo más rápido posible. Justifique su respuesta.
\end{prob}

% >>> ESPERAN PASADA CON CLAUDE CODE (pstricks) <<<

% [Original: líneas 965-987 | sin verificar | PSTRICKS - pendiente convertir a tikz | 2 variante(s) de parámetro (se descartaron 1)]
\begin{prob}
La longitud del lado marcado $x$ del triángulo de la figura (no está construído a escala) aumenta a una tasa de $0.5\ cm/s$, el lado marcado $y$ crece a una tasa de $0.3\ cm/s$ y el ángulo $\theta$ aumenta a una tasa de $0.2\ rad/s.$ Emplee la regla de la cadena para determinar la tasa a la cual el área del triángulo está cambiando en el instante $x=13\ cm,$ $y=8\ cm$ y $\theta= \dfrac{\pi}{4}.$
\textit{Sugerencia: Determine cuál es la altura del triángulo en términos de $\theta$ tomando como base a $y.$}
\begin{figure}[H]
\centering
\psset{xunit=1.0cm,yunit=1.0cm,algebraic=true,dimen=middle,dotstyle=o,dotsize=5pt 0,linewidth=1.6pt,arrowsize=3pt 2,arrowinset=0.25}
\begin{pspicture*}(-1.1,-1.64)(5.04,3.3)
\pspolygon[linewidth=2.pt](0.,0.)(3.,0.)(4.,2.)
\psline[linewidth=2.pt](0.,0.)(3.,0.)
\psline[linewidth=2.pt](3.,0.)(4.,2.)
\psline[linewidth=2.pt](4.,2.)(0.,0.)
\parametricplot{0.0}{0.4636476090008061}{0.6*cos(t)+0.|0.6*sin(t)+0.}
\begin{scriptsize}
\psdots[dotstyle=*](0.,0.)
\psdots[dotstyle=*](3.,0.)
\psdots[dotstyle=*](4.,2.)
\rput[bl](1.5,-0.32){$y$}
\rput[bl](1.84,1.3){$x$}
\rput[bl](1.02,0.2){$\theta$}
\end{scriptsize}
\end{pspicture*}
\end{figure}
\end{prob}

% [Original: líneas 1192-1214 | sin verificar | PSTRICKS - pendiente convertir a tikz | 2 variante(s) de parámetro (se descartaron 1)]
\begin{prob}
La longitud del lado marcado $x$ del triángulo de la figura aumenta a una tasa de $\dfrac{A}{8}\ cm/s$, el lado marcado $y$ crece a una tasa de $\dfrac{A}{12}\ cm/s$ y el ángulo $\theta$ aumenta a una tasa de $0.1\ rad/s.$ Emplee la regla de la cadena para determinar la tasa a la cual el área del triángulo está cambiando en el instante $x=A\ cm,$ $y=2A\ cm$ y $\theta= \dfrac{\pi}{6}.$
\textit{Sugerencia: Determine cuál es la altura del triángulo en términos de $\theta$ tomando como base a $y.$}
\begin{figure}[H]
\centering
\psset{xunit=0.8cm,yunit=0.8cm,algebraic=true,dimen=middle,dotstyle=o,dotsize=5pt 0,linewidth=2.pt,arrowsize=3pt 2,arrowinset=0.25}
\begin{pspicture*}(-0.5414253917545832,-1.6701379395708678)(6.232385046310377,3.1842198149473453)
\pspolygon[linewidth=2.pt](0.,0.)(1.8678700695702846,2.347564994457768)(5.582810856890714,2.1982317749007194)
\pscustom[linewidth=2.pt,fillstyle=solid,opacity=0.1]{
\parametricplot{0.4}{0.9}{0.4*cos(t)+0.|0.4*sin(t)+0.}
\lineto(0.,0.)\closepath}
\psline[linewidth=2.pt](0.,0.)(1.8678700695702846,2.347564994457768)
\psline[linewidth=2.pt](1.8678700695702846,2.347564994457768)(5.582810856890714,2.1982317749007194)
\psline[linewidth=2.pt](5.582810856890714,2.1982317749007194)(0.,0.)
\psdots[dotsize=4pt 0,dotstyle=*](0.,0.)
\psdots[dotstyle=*](5.582810856890714,2.1982317749007194)
\rput[bl](0.5,0.35){{$\theta$}}
\psdots[dotsize=4pt 0,dotstyle=*](1.8678700695702846,2.347564994457768)
\rput[bl](0.7588490067770819,1.3886027884036185){$x$}
\rput[bl](2.925973004329857,0.8){$y$}
\end{pspicture*}
\end{figure}
\end{prob}
